\section{Statement}

Let $\phi:\mathbb R^d\to\mathbb R\cup\{+\infty\}$ be an essentially continuous
convex function with $0<\int e^{-\phi}<\infty$. Write
\[
  Z_\phi := \int_{\mathbb R^d} e^{-\phi}\,dx, \qquad
  d\rho_\phi := Z_\phi^{-1}e^{-\phi}\,dx.
\]
For smooth strictly convex $\phi$, the \emph{Brascamp--Lieb deficit} is
\[
  \delta_{\mathrm{BL},\phi}(f)
  =
  \int_{\mathbb R^d}\bigl\langle (D^2\phi)^{-1}\nabla f,\nabla f\bigr\rangle\,d\rho_\phi
  -\Var_{\rho_\phi}(f),
\]
defined for any locally Lipschitz $f\in L^2(\rho_\phi)$. The
Brascamp--Lieb inequality \cite{BL76} asserts $\delta_{\mathrm{BL},\phi}(f)\ge 0$,
with equality cases
\[
  \mathcal O_{\mathrm{BL},\phi}
  =
  \{a\cdot \nabla\phi+b:\ a\in\mathbb R^d,\ b\in\mathbb R\}.
\]
When using $L^2(\rho_\phi)$ distances, this affine family is understood through
its members that belong to $L^2(\rho_\phi)$; for the explicit potentials
constructed below all the listed generators are in $L^2$.

\begin{theorem}\label{thm:l2-fails}
There is no finite constant $C$ such that
\begin{equation}\label{eq:would-be-stability}
  \dist_{L^2(\rho_\phi)}(f,\mathcal O_{\mathrm{BL},\phi})
  \le C\,\delta_{\mathrm{BL},\phi}(f)^{1/2}
\end{equation}
holds simultaneously for every smooth strictly convex
$\phi:\mathbb R\to\mathbb R$ with $\int e^{-\phi}<\infty$ and every locally
Lipschitz $f\in L^2(\rho_\phi)$. Consequently, no dimension-only constant
$C=C(d)$ can exist in any dimension $d\ge 1$.
\end{theorem}

\subsection*{Spectral interpretation}

In dimension one, the Brascamp--Lieb Dirichlet form is
\begin{equation}\label{eq:BL-form}
  \mathcal E_\phi(f,g)
  := \int_{\mathbb R}\frac{f'g'}{\phi''}\,d\rho_\phi,
  \qquad
  \mathcal E_\phi(f) := \mathcal E_\phi(f,f).
\end{equation}
Here the form domain means the locally absolutely continuous functions
$f\in L^2(\rho_\phi)$ for which $\mathcal E_\phi(f)<\infty$.
After closing the form, its restriction to the orthogonal complement of
constants in $L^2(\rho_\phi)$ defines a non-negative self-adjoint operator
$L_\phi$ via
$\mathcal E_\phi(f,g)=\langle L_\phi f,g\rangle_{L^2(\rho_\phi)}$, and the
Brascamp--Lieb inequality is the assertion that
$\inf\,\mathrm{spec}\bigl(L_\phi\restriction\{1\}^\perp\bigr)\ge 1$.
Theorem~\ref{thm:l2-fails} is the assertion that the spectral gap
$\inf\,\mathrm{spec}\bigl(L_\phi\restriction \mathcal O_{\mathrm{BL},\phi}^\perp\bigr)-1$
\emph{cannot} be bounded below by a positive universal constant: the
construction below produces potentials for which the Rayleigh quotient on
$\mathcal O_{\mathrm{BL},\phi}^\perp$ approaches~$1$ from above. If one also
knows compactness of the resolvent for the corresponding one-dimensional
operator, this can be phrased as collapse of the next eigenvalue.

The basic optimizer direction at the lowest level is $\phi'$ itself.

\begin{lemma}\label{lem:phi-prime-eigen}
Let $\phi:\mathbb R\to\mathbb R$ be smooth and strictly convex, with
$0<\int e^{-\phi}<\infty$. Assume in addition that $\phi''$ is bounded and
$\phi'\in L^2(\rho_\phi)$; in particular, this holds for the potentials
$\phi_S$ constructed below. Then $\int\phi'\,d\rho_\phi=0$, $\phi'$ belongs to
the form domain, and
\[
  \mathcal E_\phi(\phi',g)=\langle\phi',g\rangle_{L^2(\rho_\phi)}
  \qquad\text{for every locally Lipschitz form-domain }g.
\]
In particular $\mathcal E_\phi(\phi')=\Var_{\rho_\phi}(\phi')$, and
$\phi'$ satisfies the weak eigenvalue identity at level~$1$ against locally
Lipschitz form-domain test functions.
\end{lemma}

\begin{proof}
Integrability of $e^{-\phi}$ on $\mathbb R$ together with strict convexity
forces $\phi(x)\to+\infty$ as $|x|\to\infty$, hence $e^{-\phi(\pm R)}\to 0$
as $R\to\infty$. Moreover, by convexity $\phi'$ is monotone non-decreasing,
so $\phi'(x)$ has a (possibly infinite) limit as $x\to\pm\infty$, and there
exists $x_0>0$ with $\phi'(x)\ge 0$ for $x\ge x_0$ and $\phi'(x)\le 0$ for
$x\le -x_0$ (otherwise $\phi$ could not satisfy $\phi(\pm\infty)=+\infty$).
Thus
\[
  \int_{\mathbb R}|\phi'(x)|e^{-\phi(x)}\,dx
  = -\int_{-\infty}^{-x_0}(e^{-\phi})'\,dx
    +\int_{-x_0}^{x_0}|\phi'|e^{-\phi}\,dx
    +\int_{x_0}^{\infty}-(e^{-\phi})'\,dx <\infty,
\]
and integrating $\phi'e^{-\phi}=-(e^{-\phi})'$ over the two tails and then
letting the endpoints tend to infinity gives
$\int\phi'\,d\rho_\phi=0$.

\emph{Integration by parts for general $h$.} Let
$h\in L^2(\rho_\phi)$ be locally Lipschitz and in the form domain. Since
$\phi''$ is bounded, $h'\in L^2(\rho_\phi)$. By Cauchy--Schwarz,
$h\phi' e^{-\phi}, h e^{-\phi}, h' e^{-\phi}\in L^1(\mathbb R)$.
Choose sequences $R_k\to+\infty$ and $R_k'\to-\infty$ along which
$|h(R_k)|e^{-\phi(R_k)}+|h(R_k')|e^{-\phi(R_k')}\to0$. Integration by parts on
$[R_k',R_k]$ gives
\[
  \int_{R_k'}^{R_k}h\phi' e^{-\phi}\,dx
  = -\int_{R_k'}^{R_k}h(e^{-\phi})'\,dx
  = -\bigl[h\,e^{-\phi}\bigr]_{R_k'}^{R_k}
    +\int_{R_k'}^{R_k}h' e^{-\phi}\,dx.
\]
The boundary term vanishes along the chosen sequence; passing to the limit
yields
\begin{equation}\label{eq:ipb}
  \int_{\mathbb R} h\phi' e^{-\phi}\,dx
  = \int_{\mathbb R} h' e^{-\phi}\,dx.
\end{equation}

For any locally Lipschitz form-domain $g$, applying \eqref{eq:ipb} to $h=g$
yields
\[
  \mathcal E_\phi(\phi',g)
  =\int\frac{(\phi')'g'}{\phi''}\,d\rho_\phi
  =\int g'\,d\rho_\phi
  =\int g\,\phi'\,d\rho_\phi
  =\langle\phi',g\rangle_{L^2(\rho_\phi)}.
\]
Specializing to $g=\phi'$ gives
$\mathcal E_\phi(\phi')=\|\phi'\|_{L^2(\rho_\phi)}^2
=\Var_{\rho_\phi}(\phi')$, since $\phi'$ has $\rho_\phi$-mean zero.

For the later potentials $\phi_S$, the assumptions are immediate from the
explicit formulas: $\phi''_S$ is bounded for each fixed $S$, while on the
tails $\phi'_S(x)=\eta x\pm S$ and $e^{-\phi_S(x)}$ has at least Gaussian
decay.
\end{proof}
